\documentclass[11pt,letterpaper]{article}
% This source is intentionally restricted to pdfLaTeX.
\usepackage[T1]{fontenc}
\usepackage[utf8]{inputenc}
\usepackage{textcomp}
\usepackage{lmodern}
\usepackage[margin=0.82in,top=0.94in,bottom=0.91in,headheight=18pt]{geometry}
\usepackage{amsmath,amssymb}
\usepackage[table]{xcolor}
\definecolor{DeepNavy}{HTML}{14283F}
\definecolor{Teal}{HTML}{0A747C}
\definecolor{Gold}{HTML}{C78F36}
\definecolor{Slate}{HTML}{536477}
\definecolor{Pale}{HTML}{F4F7FA}
\definecolor{CodeEdge}{HTML}{D4DFE9}
\usepackage{microtype}
\usepackage{parskip}
\setlength{\parskip}{0.47em}
\usepackage{titlesec}
\titleformat{\section}{\normalfont\Large\bfseries\color{DeepNavy}}{\thesection}{0.62em}{}
\titlespacing*{\section}{0pt}{2.0ex plus .7ex minus .2ex}{0.75ex}
\titleformat{\subsection}{\normalfont\large\bfseries\color{Teal}}{\thesubsection}{0.6em}{}
\usepackage{longtable,booktabs,array,calc}
\renewcommand{\arraystretch}{1.22}
\AtBeginEnvironment{longtable}{\small}
\usepackage{etoolbox}
\makeatletter
\patchcmd\longtable{\par}{\if@noskipsec\mbox{}\fi\par}{}{}
\makeatother
\usepackage{listings}
\IfFileExists{upquote.sty}{\usepackage{upquote}}{}
\lstdefinestyle{suitecode}{
  basicstyle=\ttfamily\footnotesize,
  breaklines=true,breakatwhitespace=false,
  columns=fullflexible,keepspaces=true,showstringspaces=false,
  frame=single,rulecolor=\color{CodeEdge},
  backgroundcolor=\color{Pale},
  keywordstyle=\color{Teal}\bfseries,
  commentstyle=\color{Slate}\itshape,
  stringstyle=\color{DeepNavy},
  xleftmargin=5pt,xrightmargin=5pt,
  framexleftmargin=5pt,framexrightmargin=5pt,
  aboveskip=0.85em,belowskip=0.95em,
  tabsize=2
}
\lstset{style=suitecode}
\usepackage{enumitem}
\setlist[itemize]{leftmargin=1.6em,itemsep=2pt,topsep=3pt}
\setlist[enumerate]{leftmargin=1.8em,itemsep=2pt,topsep=3pt}
\providecommand{\tightlist}{\setlength{\itemsep}{0pt}\setlength{\parskip}{0pt}}
\newcommand{\passthrough}[1]{#1}
\usepackage{xurl}
\usepackage[colorlinks=true,linkcolor=Teal,urlcolor=Teal,citecolor=Teal,pdfborder={0 0 0}]{hyperref}
\urlstyle{same}
\usepackage{fancyhdr,lastpage}
\pagestyle{fancy}
\fancyhf{}
\fancyhead[L]{\footnotesize\sffamily\color{Slate} ECMA-262:2025 \textbullet{} Engineering Documentation}
\fancyhead[R]{\footnotesize\sffamily\color{Slate} 16th Edition}
\fancyfoot[L]{\footnotesize\sffamily\color{Slate} Independent technical companion}
\fancyfoot[R]{\footnotesize\sffamily\color{Slate}\thepage\ / \pageref{LastPage}}
\renewcommand{\headrulewidth}{0.35pt}
\renewcommand{\footrulewidth}{0pt}
\setlength{\emergencystretch}{3em}
\setcounter{tocdepth}{2}
\setcounter{secnumdepth}{2}
\begin{document}
\begin{titlepage}
\thispagestyle{empty}
\vspace*{2.8cm}
{\color{Gold}\rule{18mm}{3pt}}\par\vspace{18mm}
{\sffamily\Large\bfseries\color{Teal} SOFTWARE ENGINEERING DOCUMENTATION SUITE}\par
\vspace{8mm}
{\sffamily\huge\bfseries\color{DeepNavy} ECMAScript 2025 }\par
\vspace{4mm}
{\sffamily\fontsize{24}{29}\selectfont\bfseries\color{DeepNavy} Implementation Handbook }\par
\vspace{10mm}
{\large\color{Slate} Practical JavaScript engineering laboratories }\par
\vspace{16mm}
{\color{Teal}\rule{\textwidth}{0.8pt}}\par
\vfill
{\sffamily\color{Slate}\large ECMA-262:2025 \textbullet{} 16th Edition \textbullet{} June 2025}\par
\vspace{3mm}
{\sffamily\color{Slate} October 2026 \textbullet{} Document 3 of 8}\par
\vspace{3mm}
{\small\color{Slate} Original engineering companion. The uploaded ECMA-262:2025 specification remains the authoritative source.}\par
\end{titlepage}
\pagenumbering{arabic}
\tableofcontents
\clearpage
\section{Engineering workflow and technical baseline}\label{engineering-workflow-and-technical-baseline}

This is a hands-on JavaScript/ECMAScript 2025 implementation handbook. \textbf{{[}SPEC{]}} denotes behavior defined by ECMA-262:2025; \textbf{{[}HOST{]}} refers to loading, I/O, timers, web APIs, or engine selection; \textbf{{[}GUIDE{]}} denotes suggested engineering practice. Every example is original. Code using 2025 APIs requires an implementation that supports the corresponding methods; compatibility should be checked with target engines.

\textbf{Lab prerequisites:} an ECMAScript module-capable runtime, a terminal, an editor, and a test runner. Node.js commands in this guide are operational examples, not ECMA-262 requirements. Create an isolated scratch project, commit each passing lab, and do not allow shell/network side effects in untrusted code.

\section{Lab 1 --- type contracts and coercion}\label{lab-1-type-contracts-and-coercion}

\textbf{Objective.} Avoid silent data corruption at API boundaries. \textbf{{[}SPEC{]}} The \passthrough{\lstinline!Number!} type uses binary floating-point; safe integers are limited to a bounded range, while BigInt values are integers of arbitrary magnitude in the language model (6, 7, 21).

\begin{lstlisting}
export function parseQuantity(input) {
  if (typeof input !== 'string' || !/^(0|[1-9][0-9]*)$/.test(input)) {
    throw new TypeError('quantity must be canonical unsigned decimal');
  }
  const quantity = BigInt(input);
  if (quantity === 0n) throw new RangeError('quantity must be positive');
  return quantity;
}
\end{lstlisting}

\textbf{Exercise.} Reject \passthrough{\lstinline!-1!}, \passthrough{\lstinline!01!}, exponential notation, whitespace, unsafe Number coercion, and empty string. Document whether the application may accept leading zeros; the regex policy is \textbf{{[}GUIDE{]}}, not the language grammar.

\section{Lab 2 --- bindings, closures, and private state}\label{lab-2-bindings-closures-and-private-state}

\textbf{Objective.} Design APIs with predictable mutation, encapsulation, and lifetime. \textbf{{[}SPEC{]}} Lexical bindings and private class fields have distinct lookup rules; \passthrough{\lstinline!\#field!} is not equivalent to an object property named \passthrough{\lstinline!"\#field"!} (9, 15).

\begin{lstlisting}
export class RateLimiter {
  #remaining;
  constructor(tokens) {
    if (!Number.isSafeInteger(tokens) || tokens < 0) throw new RangeError('tokens');
    this.#remaining = tokens;
  }
  tryConsume() {
    if (this.#remaining === 0) return false;
    this.#remaining -= 1;
    return true;
  }
  get remaining() { return this.#remaining; }
}
\end{lstlisting}

\textbf{Exercise.} Verify external attempts to access \passthrough{\lstinline!\#remaining!} as syntax, reflection, and property keys. Add a clock-controlled replenishment policy without mixing the time source with the class's language-level invariants.

\section{Lab 3 --- prototypes, descriptors, and safe dictionary use}\label{lab-3-prototypes-descriptors-and-safe-dictionary-use}

\textbf{Objective.} Separate general object semantics from record storage. \textbf{{[}SPEC{]}} Ordinary objects have prototypes and property attributes; \passthrough{\lstinline!Object.create(null)!} creates an object without the normal inherited prototype (10, 20).

\begin{lstlisting}
export function countByKey(keys) {
  const counts = Object.create(null);
  for (const key of keys) {
    if (typeof key !== 'string') throw new TypeError('key');
    counts[key] = (counts[key] ?? 0) + 1;
  }
  return counts;
}
const result = countByKey(['toString', '__proto__', '__proto__']);
console.assert(result['__proto__'] === 2);
\end{lstlisting}

\textbf{{[}GUIDE{]}} Prefer \passthrough{\lstinline!Map!} for arbitrary keys and \passthrough{\lstinline!Object.hasOwn!} for own-property checks. Do not use \passthrough{\lstinline!for...in!} over untrusted objects without deliberate prototype handling.

\section{Lab 4 --- iterables, closing, and lazy processing}\label{lab-4-iterables-closing-and-lazy-processing}

\textbf{Objective.} Build a streaming transform that only consumes required elements. \textbf{{[}SPEC{]}} Iterators expose \passthrough{\lstinline!next()!}; iterable protocols and Iterator helper methods are specified in clauses 7.4 and 27.1. Helper operations may be lazy, while terminal operations consume elements.

\begin{lstlisting}
export function firstMatching(source, predicate) {
  return Iterator.from(source)
    .filter(predicate)
    .take(1)
    .toArray()[0];
}
console.assert(firstMatching([4, 7, 8], n => n % 2) === 7);
\end{lstlisting}

\textbf{Exercise.} Implement an iterator with \passthrough{\lstinline!next!} and \passthrough{\lstinline!return!} logging. Test whether \passthrough{\lstinline!take(1)!} closes the upstream iterator. Test thrown predicate and confirm appropriate cleanup. \textbf{{[}GUIDE{]}} Always document resource ownership when adapting file-like or network-like iterators: JS iterator closure alone does not close unrelated host resources unless the adapter implements it.

\section{Lab 5 --- immutable and set-oriented transformations}\label{lab-5-immutable-and-set-oriented-transformations}

\textbf{Objective.} Choose collection APIs with clear complexity and edge cases. \textbf{{[}SPEC{]}} \passthrough{\lstinline!Set.prototype.union!}, \passthrough{\lstinline!intersection!}, \passthrough{\lstinline!difference!}, \passthrough{\lstinline!symmetricDifference!}, \passthrough{\lstinline!isSubsetOf!}, \passthrough{\lstinline!isSupersetOf!}, and \passthrough{\lstinline!isDisjointFrom!} are defined in 24.2.

\begin{lstlisting}
export function analyzePermissions(granted, requested) {
  const g = new Set(granted);
  const r = new Set(requested);
  return {
    allowed: r.isSubsetOf(g),
    missing: r.difference(g),
    overlap: r.intersection(g),
  };
}
const check = analyzePermissions(['read', 'write'], ['read', 'admin']);
console.assert(!check.allowed && check.missing.has('admin'));
\end{lstlisting}

\textbf{{[}GUIDE{]}} This illustrates set algebra; production authorization must also check principal, resource, tenant, action, contextual policy, and revocation. Mathematical membership is not a complete authorization system.

\section{Lab 6 --- parsing, escaping, and Unicode}\label{lab-6-parsing-escaping-and-unicode}

\textbf{Objective.} Compose literal input into regular expressions without silently changing pattern semantics. \textbf{{[}SPEC{]}} \passthrough{\lstinline!RegExp.escape!} is specified in 22.2.5.1; it escapes a \emph{string to be used as pattern text}, not a serialized regular-expression object.

\begin{lstlisting}
export function containsLiteral(haystack, needle) {
  if (typeof needle !== 'string') throw new TypeError('needle');
  const pattern = new RegExp(RegExp.escape(needle), 'u');
  return pattern.test(haystack);
}
console.assert(containsLiteral('a+b', 'a+b'));
console.assert(!containsLiteral('ab', 'a+b'));
\end{lstlisting}

\textbf{Exercise.} Test strings containing \passthrough{\lstinline![!} \passthrough{\lstinline!-!} \passthrough{\lstinline!/!} newlines, surrogate pairs, and empty strings. A literal match can often use \passthrough{\lstinline!String.prototype.includes!} instead; reserve dynamic regex compilation for actual regex features. Regexp escaping does \textbf{not} prevent regex denial of service if other user-supplied pattern sections remain untrusted.

\section{Lab 7 --- async orchestration and failure propagation}\label{lab-7-async-orchestration-and-failure-propagation}

\textbf{Objective.} Normalize potentially synchronous or asynchronous callbacks without masking failures. \textbf{{[}SPEC{]}} \passthrough{\lstinline!Promise.try!} executes the callback and returns a promise that fulfills with the returned value or rejects on thrown error (27.2.4.8).

\begin{lstlisting}
export async function runTask(callback, ...args) {
  const started = Promise.try(callback, ...args);
  return await started;
}
console.assert(await runTask((x) => x + 1, 4) === 5);
try { await runTask(() => { throw new Error('fail'); }); }
catch (error) { console.assert(error.message === 'fail'); }
\end{lstlisting}

\textbf{{[}GUIDE{]}} Define cancellation and timeout policies separately; neither \passthrough{\lstinline!Promise.try!} nor \passthrough{\lstinline!Promise.race!} cancels underlying host I/O by itself. Record errors with context but never secrets.

\section{Lab 8 --- module boundaries and JSON imports}\label{lab-8-module-boundaries-and-json-imports}

\textbf{Objective.} Understand a language feature versus host interoperability. \textbf{{[}SPEC{]}} ECMAScript 2025 includes JSON module import semantics and attribute syntax in clause 16. \textbf{{[}HOST{]}} Whether \passthrough{\lstinline!./config.json!} resolves and loads depends on the host.

\begin{lstlisting}
// app.mjs, in a host supporting JSON modules and import attributes
import config from './config.json' with { type: 'json' };
export function port() {
  const candidate = config.port;
  if (!Number.isSafeInteger(candidate) || candidate < 1 || candidate > 65535)
    throw new RangeError('port');
  return candidate;
}
\end{lstlisting}

\textbf{Exercise.} Try a browser and a server runtime, documenting differences in resolution, caching, and error messages. Loading JSON is not schema validation: validate the data after loading.

\section{Lab 9 --- typed arrays and numeric fidelity}\label{lab-9-typed-arrays-and-numeric-fidelity}

\textbf{Objective.} Understand representation loss in binary data. \textbf{{[}SPEC{]}} \passthrough{\lstinline!Float16Array!}, \passthrough{\lstinline!Math.f16round!}, and DataView half-precision accessor methods are included in the 2025 edition (21, 23, 25).

\begin{lstlisting}
const h = new Float16Array(2);
h[0] = 1.5;
h[1] = 0.1;
console.assert(h[0] === 1.5);
console.assert(h[1] !== 0.1); // typically rounded to half precision
const bytes = new ArrayBuffer(2);
const view = new DataView(bytes);
view.setFloat16(0, 1.5, true);
console.assert(view.getFloat16(0, true) === 1.5);
\end{lstlisting}

\textbf{Exercise.} Document representable ranges, distinguish endianness from numerical precision, and test signed zero and non-finite values.

\section{Lab 10 --- systematic error handling}\label{lab-10-systematic-error-handling}

Build a shared error policy: unexpected programming errors propagate, user validation errors have stable types/messages, and host failures carry contextual metadata without credentials. \textbf{{[}SPEC{]}} \passthrough{\lstinline!try!}/\passthrough{\lstinline!catch!}/\passthrough{\lstinline!finally!} statements, \passthrough{\lstinline!throw!}, and Error objects are specified in clauses 14 and 20.5. \textbf{{[}GUIDE{]}} Reserve bare \passthrough{\lstinline!catch!} only when a deliberate fallback exists, and avoid swallowing abort or authorization failures.

\begin{lstlisting}
export function decodeDomainRecord(text) {
  let record;
  try { record = JSON.parse(text); }
  catch (cause) { throw new SyntaxError('invalid JSON', { cause }); }
  if (record === null || Array.isArray(record) || typeof record !== 'object')
    throw new TypeError('record must be an object');
  return record;
}
\end{lstlisting}

\section{Capstone: audited rule evaluator}\label{capstone-audited-rule-evaluator}

\textbf{Problem.} Ingest a JSON array of records containing \passthrough{\lstinline!subject!}, \passthrough{\lstinline!resource!}, and \passthrough{\lstinline!actions!}; calculate requests as sets; return an immutable summary; reject malformed types; avoid evaluation of code strings.

\textbf{Acceptance criteria:}

\begin{itemize}
\tightlist
\item[$\square$]
  Has explicit schema validation at ingress; unknown keys handled by documented policy.
\item[$\square$]
  Uses Map/Set, not inherited object keys, for arbitrary identifiers.
\item[$\square$]
  Allows no \passthrough{\lstinline!eval!}, dynamic \passthrough{\lstinline!Function!}, or arbitrary property-path execution.
\item[$\square$]
  Tests \passthrough{\lstinline!NaN!}, nullish values, empty arrays, Unicode keys, and \passthrough{\lstinline!\_\_proto\_\_!} keys.
\item[$\square$]
  Emits consistent rejection on invalid input and no unhandled promises.
\item[$\square$]
  Describes which validations are design choices and which behaviors are specified.
\end{itemize}

\section{Reviewer checklist and further reading}\label{reviewer-checklist-and-further-reading}

Reviewers should verify: input trust boundaries; numeric precision; reference/equality semantics; generator cleanup; async errors; cross-realm objects; prototype inheritance; host-dependent behavior; and feature support. Reproduce results using engine versions, not merely language-year labels.

\textbf{Primary:} ECMA-262:2025 supplied PDF, clauses 6--7, 9--10, 13--16, 20--27. \textbf{Engineering extensions:} project-specific policies above are nonnormative; applicable runtime documentation should be consulted separately.

\end{document}
